% ============================================================================
%  Inverse Kaluza-Klein Scale Symmetry
%  Document IV: Five-Dimensional Interval and Scale-Projection Mechanics
%
%  Content: 5D logarithmic scale coordinate; scale-dependent geometry;
%           complex fiber bundle; identity bundles at scale; projection
%           operator formalism; scale-dependent wavefunction; property
%           activation functions; cross-scale property map.
%
%  Cross-references: Doc II (Axioms A--J), Doc III (Kernel V4)
%  Compile: pdflatex (twice for ToC and refs)
% ============================================================================
\documentclass[11pt, a4paper, oneside]{article}

% --- Geometry ----------------------------------------------------------------
\usepackage[top=1.0in, bottom=1.0in, left=1.15in, right=1.15in,
            headheight=14pt]{geometry}

% --- Fonts / encoding --------------------------------------------------------
\usepackage[T1]{fontenc}
\usepackage{lmodern}
\usepackage{microtype}

% --- Math --------------------------------------------------------------------
\usepackage{amsmath, amssymb, amsthm, mathtools}
\usepackage{bm}
% \usepackage{physics}  % not available on this build

% --- Tables ------------------------------------------------------------------
\usepackage{booktabs}
\usepackage{array}
\usepackage{tabularx}
\usepackage{longtable}
\usepackage{multirow}

% --- Code listings -----------------------------------------------------------
\usepackage{listings}
\usepackage{xcolor}
\definecolor{codegray}{rgb}{0.5,0.5,0.5}
\definecolor{codeblue}{rgb}{0.13,0.13,0.9}
\definecolor{codegreen}{rgb}{0.0,0.5,0.0}
\definecolor{codebg}{rgb}{0.97,0.97,0.97}
\lstset{
  backgroundcolor=\color{codebg},
  basicstyle=\ttfamily\small,
  keywordstyle=\color{codeblue}\bfseries,
  commentstyle=\color{codegreen}\itshape,
  stringstyle=\color{orange},
  numberstyle=\tiny\color{codegray},
  numbers=left, numbersep=8pt,
  breaklines=true, frame=single,
  rulecolor=\color{codegray},
  captionpos=b, showstringspaces=false,
  tabsize=2
}

% --- Lists -------------------------------------------------------------------
\usepackage{enumitem}
\setlist{nosep, leftmargin=1.5em}

% --- Headers / footers -------------------------------------------------------
\usepackage{fancyhdr}
\pagestyle{fancy}
\fancyhf{}
\fancyhead[L]{\small\textit{IKK Scale Symmetry}}
\fancyhead[R]{\small\textit{Doc IV --- 5D Interval \& Projection Mechanics}}
\fancyfoot[C]{\thepage}
\renewcommand{\headrulewidth}{0.4pt}

% --- Section formatting ------------------------------------------------------
\usepackage{titlesec}
\titleformat{\section}{\Large\bfseries}{\S\thesection}{1em}{}
\titleformat{\subsection}{\large\bfseries}{\thesubsection}{1em}{}
\titleformat{\subsubsection}{\normalsize\bfseries}{\thesubsubsection}{1em}{}

% --- Hyperlinks --------------------------------------------------------------
\usepackage[colorlinks=true, linkcolor=black, citecolor=black,
            urlcolor=blue]{hyperref}

% --- Theorem-like environments -----------------------------------------------
\theoremstyle{definition}
\newtheorem{definition}{Definition}[section]
\newtheorem{axiom}{Axiom}[section]
\theoremstyle{plain}
\newtheorem{proposition}{Proposition}[section]
\newtheorem{theorem}{Theorem}[section]
\theoremstyle{remark}
\newtheorem{remark}{Remark}[section]

% --- Custom operators and macros ---------------------------------------------
%  Scale projection
\newcommand{\Proj}[2]{\Pi_{#1 \leftarrow #2}}
\newcommand{\ProjL}{\Pi_{\ell}}
%  Identity bundle
\newcommand{\Pn}{P_n}
\newcommand{\Pm}{P_m}
%  Scale level
\newcommand{\Sn}{S_n}
%  Distinguishability
\newcommand{\Dn}{D_n}
\newcommand{\Dm}{D_m}
%  Length scale
\newcommand{\Ln}{L_n}
%  Fiber
\newcommand{\Cp}{\mathbb{C}_p}
%  Activation function
\newcommand{\AQ}{A_Q(\ell)}
%  Observed / hidden states
\newcommand{\Psitrue}{|\Psi_{\mathrm{true}}\rangle}
\newcommand{\Psiobs}{|\Psi^{\mathrm{obs}}_m\rangle}
\newcommand{\Psihid}{|\Psi^{\mathrm{hid}}_{n \setminus m}\rangle}
%  5D interval
\newcommand{\dsfive}{ds_5^2}
%  Proper time
\newcommand{\dtaueff}{d\tau_{\mathrm{eff}}}

% ============================================================================
\begin{document}

\begin{titlepage}
  \centering
  \vspace*{1.5cm}
  {\large\scshape VSEPR-SIM Theoretical Division\par}
  \vspace{0.3cm}
  {\small Internal Development Document --- Continuation of Doc II\par}
  \vspace{1.2cm}
  {\Huge\bfseries Inverse Kaluza-Klein Scale Symmetry\par}
  \vspace{0.5cm}
  {\LARGE Document IV\par}
  \vspace{0.4cm}
  {\Large Five-Dimensional Interval\\[0.25em]
   and Scale-Projection Mechanics\par}
  \vspace{1.2cm}
  {\normalsize
    5D logarithmic scale coordinate\,|\,
    Scale-dependent geometry\,|\,\\[0.3em]
    Complex fiber bundle\,|\,
    Identity bundles at scale\,|\,\\[0.3em]
    Projection operator formalism\,|\,
    Scale-dependent wavefunction\,|\,\\[0.3em]
    Property activation functions\,|\,
    Cross-scale property map\par
  }
  \vspace{1.5cm}
  {\normalsize
    Prerequisites: Doc II (Axioms A--J), Doc III (Kernel V4)\\[0.4em]
    This document formalises the projection operator
    left open in Doc II Axiom~B.4 and
    the dark-scale term $\chi_w$ in Doc II Eq.~(F.5).
  \par}
  \vspace{1.5cm}
  {\large VSEPR-SIM Theoretical Division\par}
\end{titlepage}

\tableofcontents
\newpage

% ============================================================================
\section{Motivation and Starting Point}
\label{sec:motivation}
% ============================================================================

\noindent\rule{\linewidth}{0.8pt}
\textbf{Import block K} \hfill
  \textit{depends on Doc II Axioms B.2, B.3, B.4, F.5}
\noindent\rule{\linewidth}{0.4pt}

\subsection{What ordinary spacetime misses}

In special relativity an event is the four-vector
\begin{equation}
  x^\mu = (ct,\, x,\, y,\, z),
  \label{eq:4vec}
\end{equation}
and the spacetime interval is
\begin{equation}
  ds_4^2 = -c^2\,dT^2 + dx^2 + dy^2 + dz^2.
  \label{eq:SR_interval}
\end{equation}
Proper time accumulated along a worldline is
\begin{equation}
  d\tau^2 = dT^2 - \frac{1}{c^2}(dx^2 + dy^2 + dz^2).
  \label{eq:proper_time_SR}
\end{equation}
Standard result: a moving clock ages more slowly because spatial motion
consumes what would otherwise be pure temporal evolution.

\medskip
\noindent\textbf{Extension hypothesis.}
Two objects may follow identical trajectories in the four visible
coordinates yet accumulate different proper times if their internal
\emph{scale-state} evolves differently.
This demands a fifth coordinate that encodes motion through
resolution, internal phase structure, or scale-dependent access to
degrees of freedom.
It is not another place to stand; it is a record of which physical
description level the system is occupying.

\subsection{Why this addresses open variables from Doc II}

Doc II left three items undefined (see Doc II Open Variable Register,
Import Block J):
\begin{itemize}
  \item $\chi_w(w)$: the dark-scale metric correction in Eq.~(F.5).
  \item $\Pi_n$: the packing operator at each scale (Axiom B.3).
  \item $K_G(w) - K_\mathrm{EM}(w)$: the kernel difference in Eqs.~(E.3)--(E.5).
\end{itemize}
This document provides explicit geometric machinery for all three by
constructing a fifth coordinate from the physical scale variable $\ell$
and formalising the projection operator
$\Proj{m}{n}$ introduced abstractly in Doc II Axiom~B.4.

\noindent\rule{\linewidth}{0.4pt}

% ============================================================================
\section{The Five-Dimensional Interval}
\label{sec:5D_interval}
% ============================================================================

\noindent\rule{\linewidth}{0.8pt}
\textbf{Import block L} \hfill \textit{standalone; extends Eq.~(F.1) of Doc~II}
\noindent\rule{\linewidth}{0.4pt}

\subsection{Naive 5D extension}

A minimal five-dimensional interval appends one extra coordinate $u$:
\begin{equation}
  \dsfive = -c^2\,dT^2 + dr^2 + \eta_u\,du^2,
  \label{eq:5D_naive}
\end{equation}
where $dr^2 = dx^2 + dy^2 + dz^2$ and $\eta_u$ is a signature
coefficient to be determined.
This form is insufficient as stated: without specifying what $u$
measures, it says nothing.

\subsection{The scale coordinate}
\label{sec:scale_coord}

The fifth coordinate is identified with physical resolution scale:
\begin{equation}
  u = u(\ell),
  \qquad \ell \in [\ell_\mathrm{Planck},\, \ell_\mathrm{cosmo}].
  \label{eq:u_def}
\end{equation}
The natural map is logarithmic, motivated by the fact that scale
transitions span many orders of magnitude and the relevant structure
at each level is set by ratios, not absolute differences:

\begin{equation}
  \boxed{u(\ell) = \ln\!\left(\frac{\ell}{\ell_0}\right)}
  \label{eq:log_coord}
\end{equation}
where $\ell_0$ is a fixed reference length.
Consequently,
\begin{equation}
  du = \frac{d\ell}{\ell}.
  \label{eq:du}
\end{equation}

\noindent
Reference scales appearing in the theory:
\begin{center}
\begin{tabular}{lll}
  \toprule
  Level & Symbol & Approximate value \\
  \midrule
  Planck    & $\ell_\mathrm{Planck}$ & $\sim 10^{-35}$ m \\
  Nuclear   & $\ell_\mathrm{nuc}$    & $\sim 10^{-15}$ m \\
  Atomic    & $\ell_\mathrm{atom}$   & $\sim 10^{-10}$ m \\
  Molecular & $\ell_\mathrm{mol}$    & $\sim 10^{-9}$ to $10^{-7}$ m \\
  Macroscopic & $\ell_\mathrm{mac}$  & $\gtrsim 10^{-6}$ m \\
  Cosmological & $\ell_\mathrm{cosmo}$ & $\gtrsim 10^{20}$ m \\
  \bottomrule
\end{tabular}
\end{center}

\subsection{The physical 5D interval}

Substituting \eqref{eq:log_coord} into \eqref{eq:5D_naive}:

\begin{equation}
  \boxed{
    \dsfive = -c^2\,dT^2 + dr^2 + \eta_u \left(\frac{d\ell}{\ell}\right)^2
  }
  \label{eq:5D_interval}
\end{equation}

\noindent
The fifth term now explicitly measures the fractional change in resolution
scale.
A system that moves in scale (e.g.\ a particle undergoing a phase
transition between energy regimes) contributes to $ds_5^2$ through
$(d\ell/\ell)^2$, independently of its 3D spatial displacement.

\begin{remark}
  Equation~\eqref{eq:5D_interval} reduces to the standard
  4D interval~\eqref{eq:SR_interval} when $d\ell = 0$, i.e.\ when
  the description scale does not change.
  This is the correct classical limit.
\end{remark}

\subsection{Connection to the Doc II proper-time extension}

Doc II Eq.~(F.5) introduced an effective proper time
\begin{equation}
  \dtaueff = dt\,\sqrt{1 - \frac{v^2}{c^2} + \frac{2\Phi_G}{c^2} + \chi_w(w)}.
  \tag{F.5}
\end{equation}
The dark-scale correction $\chi_w(w)$ is now identified.
From~\eqref{eq:5D_interval}, the proper-time element along a 5D worldline is
\begin{equation}
  d\tau_5^2 = dT^2 - \frac{1}{c^2}\,dr^2
              - \frac{\eta_u}{c^2}\left(\frac{d\ell}{\ell}\right)^2.
  \label{eq:proper_time_5D}
\end{equation}
Comparing with (F.5):
\begin{equation}
  \chi_w \;\longleftrightarrow\; -\frac{\eta_u}{c^2}\left(\frac{d\ell}{\ell\,dT}\right)^2.
  \label{eq:chi_w_identification}
\end{equation}

The sign of $\chi_w$ is determined by $\eta_u$.
For a spacelike scale-axis ($\eta_u > 0$), motion in scale
\emph{reduces} proper time, analogously to spatial velocity.
The metric coefficient $\eta_u$ is a free parameter of the theory
until constrained by experiment; it encodes the coupling between
scale-motion and physical aging.

\noindent\rule{\linewidth}{0.4pt}

% ============================================================================
\section{Scale-Dependent Geometry}
\label{sec:scale_geometry}
% ============================================================================

\noindent\rule{\linewidth}{0.8pt}
\textbf{Import block M} \hfill \textit{standalone; prerequisite for \S\ref{sec:fibre}}
\noindent\rule{\linewidth}{0.4pt}

\subsection{Effective dimension}

Standard geometry fixes spacetime dimension at $D = 4$ for all scales.
The IKK framework replaces this with a scale-dependent effective dimension:
\begin{equation}
  D_\mathrm{eff} = D_\mathrm{eff}(\ell).
  \label{eq:Deff}
\end{equation}
The interpretation is operational: $D_\mathrm{eff}(\ell)$ counts the
number of physically accessible and dynamically active degrees of
freedom at resolution scale $\ell$.

At sub-nuclear scales, internal degrees of freedom (charge, colour,
helicity) dominate spatial ones.
At macroscopic scales, the three spatial dimensions fully activate and
internal structure averages out.
$D_\mathrm{eff}(\ell)$ interpolates this behaviour continuously.

\subsection{Scale-modulated metric}

The most general scale-dependent interval consistent with the above
writes scale-modulated coefficients on each axis:
\begin{equation}
  ds^2 = -c^2\,dT^2
         + a_1(\ell)\,dx^2
         + a_2(\ell)\,dy^2
         + a_3(\ell)\,dz^2
         + a_u(\ell)\,du^2.
  \label{eq:scale_metric}
\end{equation}
The coefficient $a_i(\ell) \in [0,1]$ measures how fully active spatial
axis $i$ is at scale $\ell$.

\begin{definition}[Effective dimension from metric coefficients]
  \begin{equation}
    D_\mathrm{eff}(\ell) = \sum_{i=1}^{3} a_i(\ell) + a_u(\ell).
    \label{eq:Deff_metric}
  \end{equation}
\end{definition}

\noindent\textbf{Limiting behaviours:}
\begin{itemize}
  \item Sub-nuclear ($\ell \lesssim 10^{-15}$ m): $a_1 \approx a_2 \gg a_3$
    may describe a system where two internal-plane directions dominate
    over the stabilised third spatial axis.
  \item Chemical / material ($\ell \sim 10^{-10}$ to $10^{-7}$ m):
    $a_1 \approx a_2 \approx a_3 \approx 1$ and ordinary 3D geometry is
    fully recovered.
  \item Cosmological ($\ell \gtrsim 10^{20}$ m): $a_u(\ell) \gg 0$,
    meaning the scale-axis becomes dynamically relevant again through
    large-scale structure formation.
\end{itemize}

\noindent\rule{\linewidth}{0.4pt}

% ============================================================================
\section{Complex Fiber Bundle}
\label{sec:fibre}
% ============================================================================

\noindent\rule{\linewidth}{0.8pt}
\textbf{Import block N} \hfill
  \textit{standalone; required for \S\ref{sec:identity_bundle}}
\noindent\rule{\linewidth}{0.4pt}

\subsection{Attaching a complex plane to every spacetime point}

A physical point is usually specified as
$p = (x, y, z, t)$.
The extended structure attaches a local complex plane at each point:
\begin{equation}
  p \;\mapsto\; (p,\, \Cp),
  \label{eq:fibre_attach}
\end{equation}
where $\Cp$ is the complex fibre above $p$.

This gives the structure of a fibre bundle
\begin{equation}
  \pi : E \to M,
  \label{eq:fibre_bundle}
\end{equation}
where $M$ is the base scale-spacetime manifold defined by the interval
\eqref{eq:5D_interval}, $E$ is the total state space, and $\pi^{-1}(p) = \Cp$.

\subsection{Local complex state}

The local state in the fibre above $p$ is written in polar form:
\begin{equation}
  \psi(p) = A(p)\,e^{i\theta(p)},
  \label{eq:local_state}
\end{equation}
where $A(p)$ is the amplitude and $\theta(p)$ is the local phase.

\begin{definition}[Wave behaviour as fibre projection]
  Wave behaviour is the observed projection of a deeper, scale-dependent
  particle identity through the complex fibres.
  The wave is not the ontologically primary object.
  It is the projection signature of an identity-bearing particle state
  onto the accessible measurement basis.
\end{definition}

\noindent
This statement is the point at which the IKK framework diverges from
``all is waves'' interpretations.
The priority is: \emph{localised identity first, wave appearance second}.

\noindent\rule{\linewidth}{0.4pt}

% ============================================================================
\section{Identity Bundles at Scale}
\label{sec:identity_bundle}
% ============================================================================

\noindent\rule{\linewidth}{0.8pt}
\textbf{Import block O} \hfill
  \textit{extends Doc II Axiom B.2 and Import Block C (packing chain)}
\noindent\rule{\linewidth}{0.4pt}

\subsection{Scale levels (Doc II recap)}

From Doc II Axiom~B.2, a scale level is the triple
\begin{equation}
  \Sn = (D_n,\, M_n,\, L_n),
  \label{eq:scale_level}
\end{equation}
where $D_n$ is the degree of freedom gained at scale $n$,
$M_n$ is the measurement basis, and $L_n$ is the characteristic
length scale.

\subsection{Identity object at scale}

A physical \emph{identity object} at scale $S_n$ is not merely a point.
It is the structured quintuple

\begin{equation}
  \boxed{P_n = \bigl(x^\mu,\; \Cp,\; D_n,\; M_n,\; L_n\bigr)}
  \label{eq:identity_bundle}
\end{equation}

The components are:
\begin{itemize}
  \item $x^\mu$: spacetime coordinates in the base manifold $M$.
  \item $\Cp$: attached complex fibre encoding local phase and amplitude.
  \item $D_n$: distinguishability content — how much identity information
    is recoverable at scale $n$.
  \item $M_n$: the measurement basis available at scale $n$.
  \item $L_n$: the length scale.
\end{itemize}

\begin{remark}
  Equation~\eqref{eq:identity_bundle} extends the Doc II Python/C++ objects
  \texttt{ScaleEntry} and \texttt{IdentityBundle} by adding the complex-fibre
  slot $\Cp$.
  The packing operator $\Pi_n$ from Doc II Axiom~B.3 maps
  $P_n \to P_{n+1}$ with compression of $D_n$.
\end{remark}

\subsection{Identity persistence across scales}

Identities at lower scales do not vanish when a higher-scale identity
assembles from them.
They remain present; the measurement system can only access their
projected contribution.

\begin{proposition}[Scale coexistence]
  For $m > n$ (coarser $>$ finer):
  \begin{equation}
    P_n \text{ exists} \;\Longrightarrow\;
    P_m = \Proj{m}{n}(P_n) \text{ coexists with } P_n.
    \label{eq:scale_coexist}
  \end{equation}
  Projection determines \emph{visibility}, not \emph{existence}.
\end{proposition}

\noindent
This avoids the emergence-as-deletion error: a molecule's formation does
not erase its constituent electrons; it hides their individual projections
behind the molecular-scale observable.

\noindent\rule{\linewidth}{0.4pt}

% ============================================================================
\section{Projection Operator Formalism}
\label{sec:projection}
% ============================================================================

\noindent\rule{\linewidth}{0.8pt}
\textbf{Import block P} \hfill
  \textit{formalises Doc II Axiom B.4; connects to Axiom B.5 residual}
\noindent\rule{\linewidth}{0.4pt}

\subsection{Definition}

\begin{definition}[Scale projection operator]
  The projection operator from scale $S_n$ to scale $S_m$ (with $m \geq n$,
  i.e.\ coarser $\geq$ finer) is
  \begin{equation}
    \Proj{m}{n} : P_n \;\longrightarrow\; P_m^{\mathrm{obs}},
    \label{eq:proj_def}
  \end{equation}
  or, in state-vector notation,
  \begin{equation}
    |\Psi_m\rangle = \Proj{m}{n}\,|\Psi_n\rangle.
    \label{eq:proj_state}
  \end{equation}
\end{definition}

\subsection{Distinguishability loss}

Projection is not copying.
It loses, compresses, averages, or re-expresses information.
Define the \emph{distinguishability} $D_n$ as the amount of
identity-resolving information at scale $S_n$.
Projection to a coarser scale satisfies
\begin{equation}
  D_m \leq D_n, \qquad m \geq n.
  \label{eq:D_loss}
\end{equation}

The lost content is
\begin{equation}
  \Delta D_{n \to m} = D_n - D_m \geq 0.
  \label{eq:D_gap}
\end{equation}

\noindent
$\Delta D_{n \to m}$ is the formal account of what appears at scale $S_m$ as:
uncertainty, phase spread, probability amplitude, quantum fuzziness,
decoherence, or unresolved wave-like behaviour.

\subsection{Decomposition identity}

The full state at scale $S_n$ decomposes exactly as:
\begin{equation}
  \boxed{
    |\Psi_n\rangle
    = \underbrace{\Proj{m}{n}|\Psi_n\rangle}_{\Psiobs}
    + \underbrace{\bigl(\mathbf{I} - \Proj{m}{n}\bigr)|\Psi_n\rangle}_{\Psihid}
  }
  \label{eq:decomp}
\end{equation}
where:
\begin{itemize}
  \item $\Psiobs = \Proj{m}{n}|\Psi_n\rangle$
    is the component visible at scale $S_m$.
  \item $\Psihid = (\mathbf{I} - \Proj{m}{n})|\Psi_n\rangle$
    is the component hidden from scale $S_m$ measurement.
\end{itemize}

\begin{remark}[Bridge between matrix and wave mechanics]
  Matrix mechanics describes the full transition structure
  $|\Psi(t)\rangle = e^{-iHt/\hbar}|\Psi(0)\rangle$.
  Wave mechanics describes the coordinate representation
  $\psi(x,t) = \langle x|\Psi(t)\rangle$.
  In the IKK framework~\eqref{eq:decomp} unifies these:
  \begin{align}
    \text{Matrix state:} &\quad
      |\Psi_n\rangle = \text{full transition structure at scale } S_n. \\
    \text{Wavefunction:} &\quad
      \psi(x,t) = \text{projected coordinate-scale representation.}
  \end{align}
  The wavefunction is not primary.
  It is the observed projection of unresolved identity content onto the
  position basis at a given scale.
\end{remark}

\subsection{Length-dependent projection}

Making the projection explicitly scale-dependent:
\begin{equation}
  \Proj{m}{n} = \Pi(L_m,\, L_n)
  \label{eq:proj_length}
\end{equation}
so that the observed state at resolution $L_m$ given a true state at $L_n$ is
\begin{equation}
  |\Psi(L_m)\rangle = \Pi(L_m, L_n)\,|\Psi(L_n)\rangle.
  \label{eq:proj_length2}
\end{equation}

For a fundamental identity, the residual hidden from all macroscopic
measurement is
\begin{equation}
  |\Psi^\perp(L)\rangle = (\mathbf{I} - \ProjL)\,\Psitrue,
  \label{eq:perp}
\end{equation}
where $\ProjL$ is the projector associated with measurement at
length scale $L$.

\subsection{Connection to the Doc II residual axiom}

Doc II Axiom~B.5 defined the residual
\begin{equation}
  \mathcal{R}_S = \mathcal{U}_S(\mathcal{I}) - \mathcal{E}_S(\mathcal{O}_S).
  \tag{B.5}
\end{equation}
This is now the projection residual of \eqref{eq:perp}:
\begin{equation}
  \mathcal{R}_S \;\equiv\; |\Psi^\perp(L_S)\rangle
               = (\mathbf{I} - \Pi_{L_S})\,\Psitrue.
  \label{eq:residual_ident}
\end{equation}
The residual is zero only when projection is injective at scale $S$,
consistent with Axiom~B.5.

\noindent\rule{\linewidth}{0.4pt}

% ============================================================================
\section{Scale-Dependent Wavefunction}
\label{sec:wavefunction}
% ============================================================================

\noindent\rule{\linewidth}{0.8pt}
\textbf{Import block Q} \hfill \textit{depends on \S\ref{sec:fibre}, \S\ref{sec:projection}}
\noindent\rule{\linewidth}{0.4pt}

\subsection{Standard form and its reinterpretation}

The standard wavefunction $\psi(x,t) = \langle x|\Psi(t)\rangle$
is the projection of a state vector onto the position basis.
In the IKK framework this is promoted to a \emph{scale-dependent} projection:

\begin{equation}
  \boxed{
    \psi(x,t;\,L) = \langle x, L\,|\,\Psi_{\mathrm{true}}\rangle
    = \Pi_{x,L}\,\Psitrue
  }
  \label{eq:wf_scale}
\end{equation}

The wavefunction is now the position projection at a specific length scale $L$.
Two observers using different resolution scales $L_1 \neq L_2$ will obtain
different wavefunctions for the same underlying identity state.

\subsection{Born rule reinterpretation}

The probability density becomes:
\begin{equation}
  |\psi(x,t;\,L)|^2
  = \text{density of projected identity-access at } (x,t)
    \text{ with resolution } L.
  \label{eq:Born}
\end{equation}

This is \emph{not} a statement that the particle ``is somewhere because
the universe enjoys gambling.''
It is a statement about how much of the deeper identity structure
is accessible through a measurement basis of resolution $L$.

\subsection{Wave-particle duality as scale-projection mismatch}

\begin{proposition}[Resolution of wave-particle duality]
  Wave-particle duality is not a fundamental ontological duality.
  It is a \emph{scale-projection mismatch}: the observer's resolution
  length $L_\mathrm{obs}$ does not fully resolve the identity length
  $L_n$ of the object being measured.
  Specifically:
  \begin{equation}
    L_\mathrm{obs} \gg L_n
    \;\Longrightarrow\;
    \Delta D_{n \to \mathrm{obs}} \gg 0
    \;\Longrightarrow\;
    \text{wave-like appearance.}
    \label{eq:duality}
  \end{equation}
  Point-like (particle) behaviour is recovered when
  $L_\mathrm{obs} \lesssim L_n$ and the projection is nearly injective.
\end{proposition}

\noindent\rule{\linewidth}{0.4pt}

% ============================================================================
\section{The Scale Ladder}
\label{sec:scale_ladder}
% ============================================================================

\noindent\rule{\linewidth}{0.8pt}
\textbf{Import block R} \hfill
  \textit{extends Doc II Scale Ladder (Import Block B)}
\noindent\rule{\linewidth}{0.4pt}

\subsection{Logarithmic scale spacing}

The scale ladder from Doc II is
\begin{equation}
  \cdots \to S_{-2} \to S_{-1} \to S_0 \to S_1 \to S_2 \to S_3 \to \cdots
  \label{eq:scale_ladder}
\end{equation}
with logarithmically spaced length scales:
\begin{equation}
  L_n = L_0\,\lambda^n,
  \label{eq:Ln}
\end{equation}
where $\lambda$ is the scale multiplier.
For decade spacing $\lambda = 10$:
\begin{equation}
  L_n = L_0 \cdot 10^n.
  \label{eq:Ln_decade}
\end{equation}

The scale coordinate of rung $n$ is
\begin{equation}
  u_n = \ln\!\left(\frac{L_n}{L_0}\right) = n\ln\lambda.
  \label{eq:un}
\end{equation}

\subsection{Projection between rungs}

The projection from rung $n$ to rung $m$ (with $m > n$, i.e.\ upward
toward coarser description) is
\begin{equation}
  |\Psi_m\rangle = \Pi(L_m, L_n)\,|\Psi_n\rangle
  = \Pi(L_0\lambda^m,\, L_0\lambda^n)\,|\Psi_n\rangle.
  \label{eq:proj_rungs}
\end{equation}

The distinguishability loss across $k = m - n$ rungs is
\begin{equation}
  \Delta D_{n \to n+k} = D_n - D_{n+k}.
  \label{eq:D_rungs}
\end{equation}
In a uniform ladder model one may write $D_n - D_{n+1} = \delta D$ (constant
per rung), giving $\Delta D_{n \to n+k} = k\,\delta D$.
The actual rung-to-rung information loss is a free parameter to be fitted
from data at each scale transition.

\noindent\rule{\linewidth}{0.4pt}

% ============================================================================
\section{Property Activation Functions}
\label{sec:activation}
% ============================================================================

\noindent\rule{\linewidth}{0.8pt}
\textbf{Import block S} \hfill \textit{depends on \S\ref{sec:scale_ladder}}
\noindent\rule{\linewidth}{0.4pt}

\subsection{Activation function definition}

A physical property $Q$ is not equally accessible at all scales.
Define the \emph{activation function}
\begin{equation}
  A_Q(\ell) \in [0,1],
  \label{eq:AQ_range}
\end{equation}
where $A_Q(\ell) = 0$ means the property is inaccessible or undefined
at scale $\ell$, and $A_Q(\ell) = 1$ means it is fully expressed.

A smooth sigmoid form consistent with the logarithmic scale coordinate is:
\begin{equation}
  \boxed{
    A_Q(\ell)
    = \frac{1}{1 + \exp\bigl[-k_Q\,(\ln\ell - \ln L_Q)\bigr]}
  }
  \label{eq:AQ}
\end{equation}
Parameters:
\begin{itemize}
  \item $L_Q$: the characteristic length at which the property activates
    (inflection point of the sigmoid in log-scale).
  \item $k_Q > 0$: sharpness of the transition.
    Large $k_Q$ approximates a step function; small $k_Q$ gives a
    smooth crossover spanning many decades.
\end{itemize}

\subsection{Observed property}

The observed value of property $Q$ at scale $\ell$ is
\begin{equation}
  Q_\mathrm{obs}(\ell) = A_Q(\ell)\;\Pi_\ell(Q_\mathrm{true}),
  \label{eq:Q_obs}
\end{equation}
where $\Pi_\ell$ is the projection operator at scale $\ell$
(see \eqref{eq:wf_scale}).

\subsection{Physical examples of activation}

\begin{itemize}
  \item \textbf{Spin.}
    Spin belongs to internal/gauge-scale identity.
    It always exists in the identity structure but activates
    magnetically only when projected to electromagnetic scales:
    \begin{equation}
      \mu_\mathrm{mag} = \Pi_{\mathrm{macro} \leftarrow \mathrm{spin}}(s),
      \quad
      A_\mathrm{spin}(\ell_\mathrm{macro}) \approx 1
      \text{ (via spin--orbit coupling)}.
      \label{eq:spin_proj}
    \end{equation}

  \item \textbf{Temperature.}
    Temperature is undefined for a single particle.
    It activates across the statistical ensemble transition:
    \begin{equation}
      k_B T \sim \frac{2}{f(\ell)}\,\langle E_\mathrm{kin}\rangle,
      \label{eq:temp_proj}
    \end{equation}
    where $f(\ell) = D_\mathrm{eff}(\ell)$ is the number of
    active degrees of freedom at scale $\ell$.
    This directly connects the activation function to
    thermodynamic behaviour: $f(\ell)$ varies from 3 (monatomic)
    to higher values as internal modes activate at larger scales.

  \item \textbf{Chemical bond geometry.}
    Bond angles are undefined at sub-nuclear scales.
    The activation length is $L_Q \sim 10^{-10}$ to $10^{-9}$ m:
    \begin{equation}
      \theta_{ijk} = \Pi_\mathrm{geometry}(\psi_i, \psi_j, \psi_k),
      \quad
      A_\mathrm{bond}(\ell) \approx 0 \;\text{ for }\;
      \ell \ll 10^{-10}\,\text{m}.
      \label{eq:bond_proj}
    \end{equation}

  \item \textbf{Inertial mass.}
    Mass is the projected resistance of identity to change of motion state:
    \begin{equation}
      m_\mathrm{obs} = \Pi_\mathrm{inertial}(P_\mathrm{true}),
      \label{eq:mass_proj}
    \end{equation}
    and for a composite body
    \begin{equation}
      M_\mathrm{obs}\,c^2
      = \sum_i m_i\,c^2 + E_\mathrm{binding} + E_\mathrm{field}.
      \label{eq:mass_composite}
    \end{equation}
    Mass is a projected property of the full identity system, not
    a simple sum of component tags.

  \item \textbf{Dark-matter-like gravitational residuals.}
    From the Doc II kernel (Eq.~E.3), the gravitational observable at
    galactic scale is
    \begin{equation}
      G_\mathrm{obs} = \Pi_\mathrm{galactic}(P_\mathrm{matter}) + \Delta G_\mathrm{scale},
      \label{eq:dark_residual}
    \end{equation}
    where $\Delta G_\mathrm{scale}$ is the residual from hidden
    scale-coordinate structure.
    The galactic-scale projection operator may expose gravitational
    behaviour that is absent at smaller scales.
    This reproduces the Doc II dark-matter kernel structure: the
    integrand $[K_G(w) - K_\mathrm{EM}(w)]$ in Eq.~(E.3) is now
    interpretable as the difference between the scale-integrated
    activation of the gravitational projection and the
    EM projection.
    \begin{equation}
      K_G(w) - K_\mathrm{EM}(w)
      \;\equiv\;
      A_G(e^w \ell_0) - A_\mathrm{EM}(e^w \ell_0).
      \label{eq:kernel_activation}
    \end{equation}
    This is a hypothesis, not a derivation.
    Numerical confirmation requires fitting $k_G$, $k_\mathrm{EM}$,
    $L_G$, $L_\mathrm{EM}$ to rotation curve data.
\end{itemize}

\noindent\rule{\linewidth}{0.4pt}

% ============================================================================
\section{Cross-Scale Property Map}
\label{sec:property_map}
% ============================================================================

\noindent\rule{\linewidth}{0.8pt}
\textbf{Import block T} \hfill
  \textit{summary table; depends on \S\ref{sec:activation}}
\noindent\rule{\linewidth}{0.4pt}

A property $Q$ is assigned the quadruple
\begin{equation}
  Q \;\mapsto\; (D_Q,\, L_Q,\, M_Q,\, \Pi_Q)
  \label{eq:prop_map}
\end{equation}
where $D_Q$ is the effective dimension required, $L_Q$ the activation
length, $M_Q$ the measurement basis, and $\Pi_Q$ the projection operator
that makes the property accessible.

\bigskip

\begin{longtable}{p{2.6cm} p{2.4cm} p{2.6cm} p{2.6cm} p{3.2cm}}
  \toprule
  \textbf{Scale regime} &
  \textbf{Length} &
  \textbf{Dominant identity} &
  \textbf{Effective $D$} &
  \textbf{Projection behaviour} \\
  \midrule
  \endhead
  Primitive / Planck &
  $10^{-35}$ m &
  Pre-geometric identity &
  Unknown / hidden &
  Fully hidden; only existence gate active \\[4pt]

  Gauge / sub-particle &
  $10^{-18}$--$10^{-15}$ m &
  Quark, lepton, gauge &
  Internal + low spatial &
  Appears as quantum number state; $|\psi\rangle$ dominated by internal channels \\[4pt]

  Nuclear &
  $10^{-15}$--$10^{-14}$ m &
  Nucleon / nucleus &
  3D onset + internal binding &
  Particle identity partially stabilised; binding energy dominant \\[4pt]

  Atomic &
  $10^{-10}$ m &
  Atom &
  3D + phase shell &
  Wavefunction dominates observable; orbital structure projects \\[4pt]

  Molecular &
  $10^{-9}$--$10^{-7}$ m &
  Molecule &
  3D structure &
  Internal phase projects into bond geometry; chirality activates \\[4pt]

  Meso\-scopic / material &
  $10^{-6}$--$10^{-3}$ m &
  Grain, domain, defect field &
  3D continuum onset &
  Many identities average into continuum fields; stress, diffusion activate \\[4pt]

  Macro\-scopic &
  $10^{-3}$--$10^{3}$ m &
  Body / object &
  Classical 3D &
  Quantum identity mostly hidden; force, torque, heat accessible \\[4pt]

  Planetary &
  $10^{6}$--$10^{8}$ m &
  Planet / system &
  3D + gravitational field &
  Internal matter projected as mass; orbital mechanics dominant \\[4pt]

  Stellar / galactic &
  $10^{9}$--$10^{21}$ m &
  Star / galaxy &
  3D + scale-gravity &
  Scale-axis residuals may reappear; rotation-curve anomalies possible \\[4pt]

  Cosmological &
  $\gtrsim 10^{21}$ m &
  Universe-scale identity &
  4D / 5D effective &
  Scale-axis dominant; expansion and large-scale structure driven by $a_u(\ell)$ \\
  \bottomrule
  \caption{Cross-scale property map: identity, dimension, and projection
    behaviour across fourteen orders of magnitude.}
  \label{tab:prop_map}
\end{longtable}

\noindent\rule{\linewidth}{0.4pt}

% ============================================================================
\section{Updated Open Variable Register}
\label{sec:open_vars}
% ============================================================================

\noindent\rule{\linewidth}{0.8pt}
\textbf{Import block U} \hfill
  \textit{updates Doc II Import Block J}
\noindent\rule{\linewidth}{0.4pt}

\begin{center}
\begin{tabularx}{\linewidth}{l l X}
  \toprule
  Variable & Status after Doc IV & Remaining requirement \\
  \midrule
  $\chi_w(w)$ &
  Identified via \eqref{eq:chi_w_identification} &
  Value of $\eta_u$ needs experimental or theoretical constraint. \\[4pt]

  $K_G(w)$, $K_\mathrm{EM}(w)$ &
  Identified as activation functions \eqref{eq:kernel_activation} &
  Parameters $k_G$, $k_\mathrm{EM}$, $L_G$, $L_\mathrm{EM}$ need fitting. \\[4pt]

  $\Pi_n$ &
  Formalised as $\Proj{m}{n}$ in \eqref{eq:proj_def}--\eqref{eq:decomp} &
  Explicit matrix form for each scale transition still needed. \\[4pt]

  $D_\mathrm{eff}(\ell)$ &
  Defined via metric coefficients \eqref{eq:Deff_metric} &
  Functional form of $a_i(\ell)$ curves needs derivation or data fit. \\[4pt]

  $\hat{H}_w$ &
  Not addressed in this document &
  Dark-scale Hamiltonian remains undefined. Required for QM predictions. \\[4pt]

  $\mathbf{T}_{AB}$ &
  Unchanged from Doc II &
  Relation-particle dynamics still undefined. \\[4pt]

  $\mathbf{W}$ &
  Unchanged from Doc II &
  Hidden-intensity weight matrix needs specific form per channel. \\[4pt]

  $k_Q$, $L_Q$ &
  New (this document) &
  Sigmoid parameters for each property require calibration against
  experimental activation energies and length-scale data. \\
  \bottomrule
\end{tabularx}
\end{center}

\noindent\rule{\linewidth}{0.4pt}

% ============================================================================
\section{Non-Technical Summary}
\label{sec:summary}
% ============================================================================

In ordinary relativity a moving clock runs slowly because spatial motion
diverts what would otherwise be pure temporal evolution.
This document extends that idea by proposing that physical systems
may also move through \emph{scale}: the same object described at
different resolutions occupies different rungs of a logarithmic
scale ladder, and this motion contributes to the proper time integral
through the fifth term $\eta_u(d\ell/\ell)^2$ in the interval
\eqref{eq:5D_interval}.

Every physical point carries an attached complex plane, making phase
and amplitude geometric features present at all scales.
What changes from one scale to the next is not whether quantum structure
exists, but how much of it can be projected into the observer's
measurement basis at that resolution.
The projection operator $\Proj{m}{n}$ maps the full identity-bearing
state at a fine scale into the observable state at a coarser one.
The complement --- the hidden component $\Psihid$ in \eqref{eq:decomp}
--- appears at the coarser scale as probability spread, phase uncertainty,
or wave-like behaviour.

In this view, wave-particle duality is not a fundamental fact about matter.
It is a scale-projection mismatch: when the resolution of the measurement
does not reach the characteristic length of the identity object, the
missing distinguishability content projects as a wave envelope.

Property activation functions~\eqref{eq:AQ} formalise the experimental
observation that certain properties --- temperature, bond geometry, stress,
dark-matter-like gravitational residuals --- are meaningless below their
characteristic activation scale but fully real above it.
The five-dimensional interval~\eqref{eq:5D_interval} provides the
geometric spine on which all of this sits.

\vfill
\begin{center}
  \rule{0.4\textwidth}{0.4pt}\\[0.5em]
  \textit{IKK Doc IV --- VSEPR-SIM Theoretical Division.
  The fifth dimension was here all along; it just needed a better address.}
\end{center}

\end{document}
